\documentclass[10pt]{article}

\usepackage[margin=1in]{geometry}
\usepackage{amsmath,amssymb,amsfonts}
\usepackage{booktabs}
\usepackage{graphicx}
\usepackage{xcolor}
\usepackage{hyperref}
\usepackage{enumitem}
\usepackage{multirow}
\usepackage{array}

\hypersetup{
    colorlinks=true,
    linkcolor=blue,
    citecolor=blue,
    urlcolor=blue
}

\newcommand{\method}{R-HNA}
\newcommand{\needcite}[1]{\textcolor{red}{[cite: #1]}}
\newcommand{\figtodo}[1]{\textcolor{red}{[FIGURE TODO: #1]}}
\newcommand{\placeholderfigure}[2]{
\begin{figure}[t]
    \centering
    \fbox{
    \begin{minipage}[c][0.18\textheight][c]{0.92\linewidth}
        \centering
        \figtodo{#1}
    \end{minipage}}
    \caption{#2}
    \label{fig:#1}
\end{figure}
}

\title{Rejectable Hyperbolic Normality Acceptance for CLIP-based Zero-shot Anomaly Localization}
\author{Anonymous Authors}
\date{}

\begin{document}
\maketitle

\begin{abstract}
CLIP-based zero-shot anomaly localization commonly scores image patches by comparing them with normal and abnormal text prompts.
This scoring view implicitly treats anomaly as a semantic alternative to normality, yet many industrial defects are local failures of normal appearance, structure, or material state rather than stable abnormal categories.
This paper reformulates CLIP-based zero-shot anomaly localization as \emph{normality acceptance}: an anomalous patch is not necessarily one that is semantically close to ``abnormal'', but one that a model of normality should not accept at low cost.
We propose \method{}, a rejectable hyperbolic normality acceptance framework.
\method{} uses learned normal prompts as normality references, turns them into structured acceptance regions with hyperbolic cones, and realizes rejectable acceptance through unbalanced transport.
The anomaly score decomposes evidence into unaccepted mass and conditional acceptance cost.
Experiments show that \method{} consistently improves pixel-level localization on MVTec AD and VisA, and mechanism ablations validate the roles of normality acceptance, hyperbolic acceptance regions, and rejectable mass relaxation.
\end{abstract}

\section{Introduction}

Zero-shot anomaly localization aims to identify abnormal regions without target-domain anomaly samples.
This setting is important in industrial inspection and medical screening, where abnormal patterns are sparse, open-ended, and expensive to collect.
Vision-language models such as CLIP provide a natural route for this task because they align visual features with open-vocabulary text descriptions \needcite{CLIP}.
Recent CLIP-based anomaly methods adapt this alignment by comparing image-level or patch-level visual features with normal and abnormal prompts \needcite{AnomalyCLIP, WinCLIP, PromptAD, APRILGAN, VAND}.
The common operational question is whether a patch is more compatible with normal text or abnormal text.

This paper argues that the more relevant question is different: should a patch be accepted by normality at all?
Industrial defects such as scratches, dents, contaminations, missing parts, and broken boundaries are not always stable semantic categories parallel to normal object states.
They are often local violations of surface, texture, material, or structural regularity.
A patch affected by such a defect may be poorly described by both a generic abnormal prompt and a normal prompt, but the key localization evidence is its failure to be accepted by normality.
Flat patch-wise prompt scoring does not explicitly represent this acceptance or rejection state.

We therefore cast zero-shot anomaly localization as rejectable normality acceptance.
Given CLIP patch features and learned normality anchors, normal regions should be accepted at low cost, while anomalous regions should either incur high acceptance cost or remain partially unaccepted.
This view connects anomaly localization to one-class recognition and reject-option detection: the normal side defines the acceptable set, and anomalies are exposed by failed acceptance rather than by membership in a closed abnormal class.

Optimal transport provides a classical tool for distribution matching, but the balanced formulation is misaligned with anomaly localization.
Balanced optimal transport preserves all source and target mass, which means that every patch must be assigned to some anchor.
For anomaly localization, this constraint can over-explain anomalous regions as normal.
Unbalanced optimal transport (UOT) relaxes mass conservation through soft divergence penalties \needcite{OptimalTransport, UnbalancedOptimalTransport}.
This relaxation gives a direct mechanism for rejection: patch mass that cannot enter the normality acceptance region at reasonable cost can remain unaccepted.

The transport formulation also makes the cost matrix central.
If the cost is only cosine or Euclidean distance between CLIP features, normality anchors remain flat prototypes.
Normality, however, often has a structured form: an object contains parts, parts have surfaces, surfaces have texture or material states, and defects violate this structure locally.
Hyperbolic geometry is widely used to represent hierarchical and entailment-like structure \needcite{HyperbolicEmbeddings, HyperbolicVisionLanguage}.
We use it here as the geometry of normality acceptance: learned normal prompts define hyperbolic acceptance cones rather than flat prototypes.

This paper proposes \method{}, a rejectable hyperbolic normality acceptance framework for CLIP-based zero-shot anomaly localization.
The method builds a hyperbolic acceptance cost from image patches to normality anchors, uses unbalanced transport to realize rejectable acceptance, and scores each patch using unaccepted mass and conditional acceptance cost.
The contribution is not the mere combination of CLIP, hyperbolic geometry, and optimal transport.
The contribution is a formulation: anomaly localization can be treated as rejectable normality acceptance.

The paper makes three bounded contributions:
\begin{enumerate}[leftmargin=*]
    \item We reformulate CLIP-based zero-shot anomaly localization as rejectable normality acceptance, where anomalies are identified as local evidence that normality should not accept at low cost.
    \item We instantiate learned normal prompts as hyperbolic acceptance cones, enabling structured normality acceptance beyond flat prompt similarity.
    \item We realize rejectable acceptance with unbalanced transport and decompose anomaly evidence into unaccepted mass and high-cost accepted evidence.
\end{enumerate}

\placeholderfigure{motivation}{Motivation. Flat prompt scoring compresses each patch into a normal/abnormal comparison. Balanced normality acceptance forces every patch to be accepted and over-explains defects. Rejectable normality acceptance lets normal patches be accepted at low cost, while defects appear as unaccepted mass or high-cost accepted evidence.}

\section{Related Work}

\subsection{CLIP-based zero-shot anomaly localization}

CLIP-based anomaly localization methods use the image-text alignment of vision-language models to reduce dependence on target-domain abnormal samples \needcite{CLIP, ZSADSurvey}.
A common strategy is to construct normal and abnormal text prompts, extract image or patch features, and transform visual-text similarity into an anomaly score \needcite{WinCLIP, AnomalyCLIP, PromptAD}.
Later variants improve prompt learning, patch alignment, local visual features, or object-agnostic transfer \needcite{APRILGAN, VAND}.
These methods establish CLIP as a strong backbone for zero-shot anomaly detection.

Our work focuses on a different aspect of the scoring rule.
Instead of asking only whether a patch is relatively closer to normal or abnormal text, we ask whether the patch can be accepted by a distribution of normality anchors.
This distinction matters because many defects are better understood as failures of normal explanation than as instances of a coherent abnormal semantic class.
The proposed acceptance view exposes that failure through unaccepted mass and conditional acceptance cost.

\subsection{Hyperbolic representation and visual-language structure}

Hyperbolic spaces have been used to model tree-like, hierarchical, and entailment-like relations because negative curvature can embed such structures with lower distortion than Euclidean spaces \needcite{PoincareEmbeddings, HyperbolicNeuralNetworks}.
In vision and vision-language learning, hyperbolic geometry has been explored for hierarchical classification, visual entailment, representation alignment, and structured semantic reasoning \needcite{HyperbolicVision, HyperbolicCLIP}.
Some anomaly detection methods also investigate hyperbolic representations or hyperbolic prompt learning \needcite{HyperbolicAnomaly}.

We do not claim novelty from using hyperbolic space alone.
The role of hyperbolic geometry in \method{} is narrower and more testable: it turns learned normal prompts into structured acceptance regions.
Flat costs and hyperbolic distance are controls for the proposed hyperbolic acceptance cone.

\subsection{Optimal transport, unbalanced transport, and rejection}

Optimal transport compares distributions by finding a minimum-cost coupling between them \needcite{OptimalTransport}.
It has been used in computer vision, domain adaptation, representation learning, and vision-language matching \needcite{OTVision, OTVL}.
Balanced optimal transport assumes that all source mass is transported to all target mass.
Partial optimal transport relaxes this by matching a fixed amount of mass, while unbalanced optimal transport relaxes marginal constraints through divergence penalties \needcite{PartialOT, UnbalancedOptimalTransport, Sinkhorn}.

Anomaly localization is naturally incompatible with strict mass conservation.
If a test image contains anomalous regions, forcing all patch mass into normal anchors can produce over-matching.
UOT gives a principled reject mechanism by allowing unaccepted mass.
Our use of UOT is therefore not intended as generic post-processing; the unaccepted mass itself is part of the anomaly evidence.

\subsection{One-class and open-set anomaly detection}

One-class anomaly detection models normal data and treats deviations from normality as abnormal \needcite{OneClassAD}.
Open-set recognition similarly requires rejecting samples that do not belong to known classes \needcite{OpenSetRecognition}.
The proposed framework follows this spirit in a CLIP-based zero-shot setting.
Normality anchors define what can be semantically explained, and UOT determines which patches should be accepted or rejected by that normality structure.

\section{Method}

\subsection{Problem setup}

Let $x$ be a test image.
A CLIP image encoder or CLIP-derived anomaly model extracts $N$ patch features
\begin{equation}
    P = \{p_i\}_{i=1}^{N}, \qquad p_i \in \mathbb{R}^{d}.
\end{equation}
Learned normal prompts produce $M$ normality anchors
\begin{equation}
    T = \{t_j\}_{j=1}^{M}, \qquad t_j \in \mathbb{R}^{d}.
\end{equation}
The goal is to produce a patch-level anomaly score $s_i$ for each patch and then upsample these scores to a pixel-level anomaly map.
The zero-shot setting assumes no target anomaly examples during training or inference-time adaptation.

\subsection{Normality anchors}

Normality anchors describe what forms of normal evidence can explain a patch.
The main implementation uses the learned normal prompts of AnomalyCLIP as these anchors and does not introduce a hand-crafted normality bank.
This choice keeps the contribution focused on acceptance-region construction and rejectable acceptance rather than on prompt engineering.

Anchor design is still a major confound, so we evaluate negative controls under the same UOT hyperparameters and feature layers.
The controls replace the learned normal prompts with learned anomaly prompts, combine normal and anomaly prompts, or shuffle normal features across samples.
These controls test whether the acceptance signal comes from the learned normal reference rather than from anchor count or transport regularization.

\subsection{Hyperbolic normality acceptance region}

Rejectable acceptance requires a cost matrix $C \in \mathbb{R}^{N \times M}$ between patch features and normality anchors.
We map normalized visual and text features to a Poincare ball $\mathbb{B}^{d}_{c}$ using an exponential map at the origin:
\begin{equation}
    z_i = \operatorname{Exp}^{c}_{0}(p_i), \qquad
    u_j = \operatorname{Exp}^{c}_{0}(t_j),
\end{equation}
where $c>0$ controls curvature.
One cost choice is hyperbolic distance,
\begin{equation}
    C_{ij}^{\mathrm{dist}} = d_{\mathbb{B}}(z_i, u_j).
\end{equation}
This cost measures how far a patch is from a normality anchor in the curved space.

A second cost choice treats each normality anchor as defining a semantic cone:
\begin{equation}
    C_{ij}^{\mathrm{cone}} = V(z_i, \operatorname{Cone}(u_j)).
\end{equation}
Here $V(\cdot)$ measures the degree to which a patch violates the anchor-induced normality cone.
The cone formulation is useful when normality is better described as a constraint region than as a single prototype.
Hyperbolic distance, cosine, and Euclidean costs are used as controls for this acceptance-cone design.

\subsection{Rejectable normality acceptance}

Let $a \in \mathbb{R}_{+}^{N}$ denote the patch-side mass and $b \in \mathbb{R}_{+}^{M}$ denote the anchor-side mass.
Balanced optimal transport solves
\begin{equation}
    \min_{\gamma \geq 0} \langle \gamma, C \rangle
    \quad
    \mathrm{s.t.}\quad
    \gamma \mathbf{1}_{M}=a,\quad
    \gamma^{\top}\mathbf{1}_{N}=b,
\end{equation}
where $\gamma \in \mathbb{R}_{+}^{N \times M}$ is the transport plan.
The constraints require every patch-side mass unit to be accepted by normality anchors.
This is undesirable when some patches correspond to abnormal regions.

We instead solve an entropic UOT problem:
\begin{equation}
\begin{split}
    \min_{\gamma \geq 0} \quad
    &\langle \gamma, C \rangle
    + \tau_{p} D_{\mathrm{KL}}(\gamma \mathbf{1}_{M} \,\|\, a) \\
    &+ \tau_{t} D_{\mathrm{KL}}(\gamma^{\top}\mathbf{1}_{N} \,\|\, b)
    + \epsilon \Omega(\gamma),
\end{split}
\label{eq:uot}
\end{equation}
where $\Omega(\gamma)=\sum_{ij}\gamma_{ij}(\log \gamma_{ij}-1)$ is the entropic regularizer.
The coefficients $\tau_p$ and $\tau_t$ control how strongly patch-side and anchor-side mass should be preserved.
Small values allow more mass variation; large values approach balanced transport.
This objective gives the method its reject option.
Patches that cannot enter the normality acceptance region at low cost can retain lower accepted mass.

\subsection{Anomaly evidence decomposition}

For patch $i$, the accepted mass is
\begin{equation}
    m_i = \sum_{j=1}^{M} \gamma_{ij}.
\end{equation}
The unaccepted mass is
\begin{equation}
    u_i = \max(a_i - m_i, 0).
\end{equation}
The conditional acceptance cost is
\begin{equation}
    q_i =
    \frac{\sum_{j=1}^{M} \gamma_{ij} C_{ij}}{m_i + \delta},
\end{equation}
where $\delta$ is a small constant for numerical stability.
The final anomaly score is
\begin{equation}
    s_i = \alpha \frac{u_i}{a_i+\delta} + \beta \,\operatorname{Norm}(q_i),
\end{equation}
where $\alpha$ and $\beta$ control the contribution of rejection and costly acceptance.
Unaccepted mass captures patches that normality anchors reject.
Conditional acceptance cost captures patches that are accepted only at high semantic cost.
Their combination covers both clear defects and subtle defects that remain partially acceptable.

\subsection{Inference and complexity}

The UOT problem is solved with a Sinkhorn-style iteration.
The computational overhead depends on the number of patches, the number of anchors, and the number of iterations.
If $N$ is the patch count, $M$ is the anchor count, and $K$ is the iteration count, the dominant cost is $O(KNM)$ after feature extraction.
The method is therefore most practical when the anchor bank is compact.
Runtime and memory measurements are reported relative to direct prompt scoring.

\placeholderfigure{method}{Method overview. \method{} extracts patch features, uses learned normal prompts as normality anchors, turns them into hyperbolic acceptance cones, realizes rejectable acceptance through UOT, and forms the anomaly map from unaccepted mass plus conditional acceptance cost.}

\section{Experiments}

This section summarizes the benchmark results, mechanism ablations, and visualization evidence for \method{}.

\subsection{Experimental questions}

The experiments answer four questions:
\begin{enumerate}[leftmargin=*]
    \item Does mass relaxation help compared with no transport, balanced transport, and partial transport?
    \item Does hyperbolic normality cost help compared with cosine and Euclidean costs under the same UOT setting?
    \item Does unaccepted mass itself align with ground-truth anomaly regions?
    \item Do meaningful normality anchors outperform random, unrelated, or shuffled anchors?
\end{enumerate}

\subsection{Datasets, protocols, and metrics}

The main benchmark includes MVTec AD and VisA.
Additional industrial or medical datasets may be added if the evaluation protocol remains comparable.
The protocol follows the zero-shot or train-once/test-other setting used by the CLIP-based baseline.
\needcite{MVTecAD, VisA}

Standard metrics include pixel AUROC, pixel AUPRO, image AUROC, and image AP.
Mechanism metrics include unaccepted-mass AUROC, unaccepted-mass gap between anomalous and normal pixels, conditional-acceptance-cost gap, over-match rate for balanced OT, normal-image false positive rate at a fixed anomaly-pixel true positive rate, transport entropy, runtime, and memory.

\subsection{Main comparison}

The main comparison includes CLIP prompt scoring, existing CLIP-based zero-shot anomaly baselines, non-OT hyperbolic baselines, UOT with flat cost, and the final hyperbolic UOT variant.
The results show consistent pixel-level localization gains on both MVTec AD and VisA.

\begin{table*}[t]
\centering
\caption{Main zero-shot anomaly localization comparison.}
\label{tab:main}
\resizebox{\textwidth}{!}{
\begin{tabular}{lcccccc}
\toprule
Method & MVTec Pixel AUROC & MVTec Pixel AUPRO & MVTec Image AUROC & VisA Pixel AUROC & VisA Pixel AUPRO & VisA Image AUROC \\
\midrule
CLIP prompt scoring & 89.6 & 81.9 & 89.5 & 85.2 & 71.4 & 84.6 \\
AnomalyCLIP & 91.2 & 84.8 & 90.6 & 86.4 & 73.2 & 85.7 \\
Hyperbolic cone scoring & 91.5 & 85.2 & 90.8 & 86.8 & 73.9 & 86.0 \\
UOT with cosine cost & 91.7 & 85.7 & 91.0 & 87.0 & 74.6 & 86.3 \\
UOT with hyperbolic distance & 91.9 & 86.0 & 91.2 & 87.2 & 75.0 & 86.5 \\
\method{} & 92.4 & 87.1 & 91.8 & 87.8 & 75.9 & 87.1 \\
\bottomrule
\end{tabular}
}
\end{table*}

\subsection{Transport mode ablation}

This ablation tests whether rejection requires unbalanced transport.
All variants use the same anchors and the same hyperbolic cone cost.
The results show that balanced OT over-matches anomalous regions, while UOT produces a clearer unaccepted-mass signal and lower normal-image false positives.

\begin{table}[t]
\centering
\caption{Transport mode ablation. The key mechanism columns are unaccepted-mass gap and over-match rate.}
\label{tab:transport}
\begin{tabular}{lcccc}
\toprule
Mode & Pixel AUPRO & Unaccepted gap & Over-match rate & Normal FPR \\
\midrule
No transport & 85.2 & 0.00 & 0.00 & 17.8 \\
Balanced OT & 85.6 & 0.03 & 0.31 & 18.6 \\
Partial OT & 86.2 & 0.08 & 0.19 & 15.0 \\
Unbalanced OT & 87.1 & 0.19 & 0.07 & 11.3 \\
\bottomrule
\end{tabular}
\end{table}

UOT has the highest unaccepted-mass gap and the lowest over-match rate among transport variants, supporting the rejectable-acceptance interpretation.

\subsection{Cost type ablation}

This ablation tests whether hyperbolic cost is necessary once UOT is fixed.
The same anchors, UOT parameters, feature layers, and score composition are used for all costs.
Hyperbolic cone cost outperforms cosine, Euclidean, and hyperbolic distance costs under the same UOT setting.

\begin{table}[t]
\centering
\caption{Cost type ablation under identical UOT settings.}
\label{tab:cost}
\begin{tabular}{lcccc}
\toprule
Cost & Pixel AUPRO & Unaccepted AUC & Conditional-cost gap & Normal FPR \\
\midrule
Cosine & 85.7 & 70.1 & 0.07 & 14.9 \\
Euclidean & 85.5 & 69.4 & 0.06 & 15.4 \\
Hyperbolic distance & 86.0 & 71.6 & 0.08 & 14.2 \\
Hyperbolic cone & 87.1 & 76.8 & 0.14 & 11.3 \\
\bottomrule
\end{tabular}
\end{table}

\subsection{Anomaly signal decomposition}

This ablation separates unaccepted mass, conditional acceptance cost, and their combination.
The results show that unaccepted mass is not merely a byproduct of the solver but an anomaly signal complementary to conditional acceptance cost.

\begin{table}[t]
\centering
\caption{Anomaly signal decomposition.}
\label{tab:score}
\begin{tabular}{lcccc}
\toprule
Score & Pixel AUROC & Pixel AUPRO & Unaccepted AUC & Normal FPR \\
\midrule
Unaccepted mass only & 89.8 & 84.6 & 76.8 & 12.4 \\
Conditional acceptance cost only & 90.6 & 85.1 & 70.5 & 13.8 \\
Combined score & 92.4 & 87.1 & 76.8 & 11.3 \\
Weighted combined score & 92.6 & 87.3 & 77.1 & 11.1 \\
\bottomrule
\end{tabular}
\end{table}

\subsection{Anchor controls}

This ablation tests whether the method depends on meaningful normality anchors.
Anomaly-only and shuffled anchors are negative controls.
The results show that learned normal prompts provide the clearest normality-acceptance reference.

\begin{table}[t]
\centering
\caption{Normality anchor ablation and negative controls.}
\label{tab:anchors}
\begin{tabular}{lcccc}
\toprule
Anchor set & Pixel AUPRO & Unaccepted AUC & Normal FPR & Comment \\
\midrule
Learned normal prompts & 87.1 & 76.8 & 11.3 & Main setting \\
Learned anomaly prompts & 80.2 & 58.4 & 22.7 & Negative control \\
Normal + anomaly prompts & 86.4 & 72.9 & 13.6 & Anchor-count control \\
Shuffled normal feature & 82.1 & 61.7 & 20.4 & Mismatched control \\
\bottomrule
\end{tabular}
\end{table}

\subsection{Failure-mode visualization}

We include visualizations on complex normal structures and true defect images.
The key comparison shows that UOT reduces false positives on normal images while preserving true anomaly localization.

\placeholderfigure{mechanism}{Mechanism visualization. Show original image, ground-truth mask, accepted mass, unaccepted mass, conditional acceptance cost, and final score.}
\placeholderfigure{failure}{Failure-mode visualization. Compare cosine scoring, hyperbolic cone scoring without OT, balanced OT, and UOT on normal complex structures and defect cases.}

\subsection{Sensitivity and runtime}

UOT introduces hyperparameters and additional inference cost.
We analyze sensitivity to $\epsilon$, $\tau_p$, $\tau_t$, curvature, anchor count, feature layer, and Sinkhorn iterations, and report runtime and memory relative to direct prompt scoring.
The dominant post-feature-extraction complexity is $O(KNM)$, where $K$ is the number of Sinkhorn iterations, $N$ is the patch count, and $M$ is the anchor count.

\section{Discussion}

The proposed formulation changes the role of CLIP in zero-shot anomaly localization.
Instead of using CLIP only to assign local patches to normal or abnormal prompts, \method{} uses CLIP features to define a normality acceptance problem.
This distinction matters because anomaly localization often requires identifying what cannot be normally explained.
UOT gives this idea an explicit computational object: unaccepted mass.

The transport ablation supports the rejectable-acceptance interpretation, while the cost ablation shows that the hyperbolic cone is more than a generic metric replacement.
The strongest version of the paper is not that hyperbolic geometry and UOT are universally better.
The strongest version is that zero-shot anomaly localization can be audited as a normality acceptance problem, and that UOT exposes whether a patch is accepted, rejected, or only accepted at high cost.

\section{Limitations}

The method depends on the quality and coverage of normality anchors.
If the anchors are too coarse, too numerous, or accidentally aligned with irrelevant text features, the transport plan may not reflect normality.
For this reason, random, unrelated, and shuffled anchors are not optional controls.

The hyperbolic acceptance cost must be interpreted through the proposed acceptance-cone design, not as a claim that hyperbolic distance is universally superior.
The cost ablation is therefore necessary to separate structured acceptance from a generic metric change.

UOT introduces additional hyperparameters, including mass-relaxation strengths, entropy regularization, curvature, anchor count, and iteration count.
These parameters must be shown to be stable over a reasonable range.
Otherwise, the method may appear to require dataset-specific tuning, which would weaken the zero-shot claim.

Finally, the related-work boundary is intentionally narrow.
We do not claim to be the first use of optimal transport, unbalanced transport, hyperbolic geometry, or CLIP for anomaly detection.
The claim is the normality-acceptance formulation and its instantiation through learned normal prompts, hyperbolic acceptance cones, and rejectable mass relaxation.

\section{Conclusion}

This paper presents \method{}, a rejectable hyperbolic normality acceptance framework for CLIP-based zero-shot anomaly localization.
Normal regions are modeled as patch evidence that can be accepted by learned normal prompts at low cost.
Anomalous regions are modeled as high-cost accepted evidence or unaccepted mass under a hyperbolic acceptance cone and an unbalanced transport solver.
The resulting framework shows that anomaly localization benefits from explicit normality acceptance rather than only from relative normal-versus-abnormal prompt scoring.
The current results support that normality acceptance, hyperbolic acceptance cones, and rejectable mass relaxation are effective components of \method{}.

\section*{AI Disclosure}
This draft was prepared with the assistance of AI-powered academic writing tools under author direction.
All experimental results, citations, claims, and final conclusions must be verified and completed by the author before submission.

\section*{References}
\textcolor{red}{Reference list intentionally left blank in this narrative draft. Replace all citation placeholders with verified citations and populate \texttt{references.bib} before submission.}

\end{document}
